\documentclass[11pt]{article}
\usepackage[margin=0.5in]{geometry}

\newcommand{\duedate}{}
\newcommand{\assignment}{Summary of Weimer et al. (2023)}

\newcommand{\name}{Aksel Kretsinger-Walters, adk2164}
\newcommand{\email}{adk2164@columbia.edu}

\makeatletter
\def\input@path{{../}{../../}{../../../}}
\makeatother
\input{pset_template_6873.tex}

\begin{document}
\psetheader

\section*{Weimer et al. (2023): A Causal Approach for Detecting Team-Level Momentum in NBA Games}

\subsection*{1. Research Question}

Does \textbf{team-level momentum} exist in basketball?

Momentum hypothesis:
\[
		P(\text{score next possession} \mid \text{recent scoring run})
		>
		P(\text{score next possession})
\]

Traditional skepticism:
\begin{itemize}
		\item Many studies treat scoring as independent events.
		\item Momentum often attributed to cognitive bias (``hot hand illusion'').
\end{itemize}

This paper tests whether \textbf{timeouts interrupt momentum}.
If pauses reduce scoring by the momentum team, momentum must exist.

\bigskip
Key idea:

\[
		\text{If momentum does not exist} \Rightarrow \text{timeouts should not affect scoring}.
\]

\subsection*{2. Causal Identification Strategy}

Problem:
\begin{itemize}
		\item Coaches call timeouts strategically.
		\item Strategic timeouts introduce \textbf{selection bias}.
\end{itemize}

Solution: Use \textbf{TV timeouts}.

TV timeouts are:
\[
		\text{exogenous stoppages of play}
\]

They occur at predetermined times in the NBA rulebook and are unrelated to momentum.

Thus:

\[
		\text{TV timeout} \approx \text{random assignment of pause in play}
\]

This approximates a randomized experiment.

\subsection*{3. Matching Framework}

Goal: Compare runs that are identical except for the presence of a TV timeout.

Matching variables:
\begin{itemize}
		\item run length
		\item substitutions by momentum team
		\item substitutions by opposing team
\end{itemize}

Propensity score matching:

\[
		Pr(Treatment_i = 1 \mid X_i)
\]

where

\[
		X_i = (\text{run size}, \text{substitutions})
\]

Each treated run (interrupted by TV timeout) is matched with a similar untreated run.

\subsection*{4. Operational Definition of Momentum}

Momentum indicator:

\[
		m_i =
		\begin{cases}
				1 & \text{if } p_u \ge 6 \\
				0 & \text{otherwise}
		\end{cases}
\]

where

\[
		p_u = \text{unanswered points}.
\]

Thus momentum begins when:

\[
		\text{run} \ge 6 \text{ points}.
\]

Robustness checks also test:

\[
		10 \text{ point runs}, \quad 15 \text{ point runs}.
\]

Outcome variable:

\[
		Y_i = \text{points scored by the momentum team in the next 3 minutes}.
\]

\subsection*{5. Statistical Model}

Outcome modeled with Poisson regression:

\[
		Y_i \sim \text{Poisson}(\lambda_i)
\]

\[
		\log(\lambda_i) =
		\beta_0
		+ \beta_1 \text{TVTimeout}_i
		+ \beta_2 \text{RunSize}_i
		+ \beta_3 \text{SubMomentum}_i
		+ \beta_4 \text{SubOpponent}_i
\]

Key coefficient:

\[
		\beta_1 < 0
\]

indicates that timeouts reduce scoring during momentum.

\subsection*{6. Main Result}

TV timeouts reduce scoring by the momentum team:

\[
		e^{-0.119} - 1 \approx -11.2\%
\]

Interpretation:

\[
		\text{Momentum team scores } \approx 11\% \text{ fewer points after TV timeout}.
\]

Average effect:

\[
		\approx 0.56 \text{ fewer points over the next 3 minutes}.
\]

Stronger effects appear for longer runs.

\subsection*{7. Mechanisms}

Two mechanisms explain the effect:

\textbf{1. Pause in play}

Stopping gameplay interrupts momentum.

\textbf{2. Strategy changes}

Substitutions or tactical adjustments.

Evidence:

\begin{itemize}
		\item Timeouts with no substitutions still reduce scoring.
		\item Strategy explains only about 40\% of the effect.
\end{itemize}

Thus most of the effect comes from the \textbf{stoppage itself}.

\subsection*{8. Robustness Checks}

Several tests confirm results:

\begin{itemize}
		\item different run thresholds (6, 10, 15 points)
		\item alternative time windows for scoring
		\item negative binomial models
		\item Monte Carlo simulation with no momentum
\end{itemize}

Simulation result:

\[
		\text{If momentum } = 0
		\Rightarrow \text{timeouts have no effect}.
\]

Thus the empirical effect is not a statistical artifact.

\subsection*{9. Interpretation}

Momentum behaves like a \textbf{state variable}:

\[
		\text{performance}_{t+1}
		=
		f(\text{performance}_t)
\]

Pausing play disrupts that state.

This implies:

\begin{itemize}
		\item momentum exists at the \textbf{team level}
		\item interruptions reduce performance
\end{itemize}

\subsection*{10. Implications}

Strategic insights:

\begin{itemize}
		\item Timeouts help stop opponent momentum.
		\item Momentum teams should avoid substitutions.
\end{itemize}

Methodological contribution:

\begin{itemize}
		\item Demonstrates how \textbf{natural experiments} can detect causal effects in sports.
		\item Shows value of \textbf{matching methods in observational data}.
\end{itemize}

\bigskip

\textbf{Key Insight}

Momentum in basketball is not purely psychological.

\[
		\text{Performance is serially correlated at the team level.}
\]

Interrupting play breaks that correlation.

\end{document}
